\section{当领域无法快速验证：慢环境、reward model 与分解}

\videofigure{figures/self-improvement-directions.jpg}{未来自我改进需要更强的跨域泛化、robust/meta verification，以及自动选择 train-time scaling 数据。}{00:33:34--00:40:08}

数学和代码之所以率先成功，是因为 reward 快且明确。芯片设计仿真、科学实验、机器人和创意任务可能需要数小时、数天或主观判断，无法在 RL 中执行成千上万次。

\begin{importantbox}{不可快速验证不等于完全不可验证}
可以把昂贵真实评估转为稀疏 gold signal，再训练 surrogate reward model；也可把大任务分解为编译、局部仿真、性能 profile 等较快子验证器。但代理信号必须定期用真实环境校准。
\end{importantbox}

一种近似方式是学习 reward model $\hat R_\phi(x,y)$ 预测昂贵环境回报 $R(x,y)$：

$$
\phi^*=\arg\min_\phi
\mathbb E_{(x,y,R)\sim \mathcal D}
\left[\hat R_\phi(x,y)-R(x,y)\right]^2.
$$

随后 RL 在 $\hat R_\phi$ 上快速迭代，但策略会主动寻找 reward model 的盲点，因此需要 active sampling 把高不确定样本送回真实仿真或专家。

\begin{warningbox}{Surrogate reward 的泛化范围决定安全上限}
如果策略走出训练数据分布，reward model 可能高分奖励物理上无效或危险的设计。模型越会优化，越会放大 evaluator 的小误差。
\end{warningbox}

\begin{knowledgebox}{任务分解是能力桥梁，不是最终答案}
当前模型无法优化完整芯片时，可以先做 kernel、局部模块和 profile 分析。分解让现有能力参与复杂流程，但最终仍要验证子任务组合后是否满足全局目标。
\end{knowledgebox}

\subsection{本章小结}

在慢验证和主观领域，自我改进需要 surrogate reward、分层任务与稀疏真实反馈。关键挑战从“有没有 reward”变成“代理 reward 在策略优化后是否仍可信”。

